% =============================================================
% §6 讨论
% =============================================================
\section{讨论}
\label{sec:discuss}

\subsection{对 AIGC 调度研究的启示}
\label{sec:implications}

我们的实验支持三条超越 \S\ref{sec:eval} 所比较的具体调度器的启示。

\textbf{仿真器物理的影响力远超算法精巧度。}本研究中最大的效应---关闭连续批处理造成的 $-83\%$ \slo{} 坍缩（表~\ref{tab:ablation}）---比我们测得的任何奖励设计或状态工程选择都大一个数量级。研究精力与其投入网络架构或奖励公式，不如投入忠实的物理仿真。

\textbf{预训练是性价比最高的必备投资。}关闭预训练造成 $38\%$ \slo{} 跌幅（$p<0.001$），但只需 $20$ 个 episode（约 10--15~s 墙钟）即可。预训练的意义不在样本效率，而在于\textbf{跳出局部最优区域}：早期配置都陷入"always-cloud"吸引子，直到我们将充足预训练与 cloud-overuse 惩罚联合引入。

\textbf{勿过度归因于奖励整形。}一旦连续批处理与预训练就位，没有任何单一 AIGC 奖励项（乃至全部联合）能显著改善 \slo{}；权重敏感性研究（\S\ref{sec:weight-sens}）进一步表明，全部 7 项的\emph{相对权重}可从延迟优先切换到 AIGC 优先而无可测量的影响。先前声称某具体奖励项贡献的工作，可能忽略了充足表达力的 PPO 主干内部这种\textbf{可替换性}：缺少表~\ref{tab:ablation} 这类联合消融时，单项奖励的正向结果\emph{值得怀疑}。

总体上，这些启示提示 AIGC 调度研究议程应当从"造更好的奖励"转向\textbf{"造更好的仿真器与训练流程"}。

\subsection{局限与有效性威胁}
\label{sec:limitations}

\textbf{为何用仿真，及其保真度。}我们刻意选择在仿真中评估。本文的核心论断是\emph{比较性、统计性}的---一个必须在 30 种配置、每种 $N{=}30$ 个种子（共 900 次运行）上都成立的 Pareto 占优关系。在真实硬件上复现这一扫描，需要同时获得 A100、T4、Jetson 级节点并施加可控可复现的到达过程---而这恰是仿真器所隔离的场景：全部 11 种调度器跑在完全相同的负载与种子上，不受多租户、热降频、驱动漂移等混杂因素干扰。我们将仿真器常数（冷加载时间、批处理开销、显存带宽下限）标定至 vLLM 与 TensorRT-LLM 的公开实测值，使绝对量级而非仅相对量级保持真实。

尽管如此，仿真器有两处简化：将连续批处理建模为\emph{准入式}批处理而非 vLLM 的迭代级动态再批处理（在某些工况下可能低估真实批处理收益 10--20\%），并将显存带宽下限视为常数而非随 batch size 变化。因此我们把绝对消融量级（如 \texttt{no\_batching} 的 $-83\%$）理解为仿真器行为而非 vLLM 精确预测；而承载本文论断的效应\emph{排序}与\emph{显著性}对这些常数稳健。\textbf{在配备生产级推理引擎的物理云--边测试床上验证该 Pareto 关系，是最重要的未来工作}；我们开源的仿真器正是为让这一对比直接可行而设计：其 \texttt{scheduler.schedule()} 接口与真实派发器需实现的接口完全一致。

\textbf{工作负载与硬件覆盖。}我们的负载基于对齐 Azure 2023 的对数正态长度合成数据---未回放生产 trace，也未覆盖突发到达---硬件涵盖 A100/T4/Jetson 层级，但未含 H100 或定制加速器。所建模的定性物理特性架构无关；常数需对新硬件重新标定。

\textbf{统计功效与算法范围。}$N=30$ 次运行可检测中到大效应（Cohen's $d \ge 0.5$）；消融的零结果意味着"未能检测到大于 5\% 的 \slo{} 差异"，而非"不存在差异"。PPO、A3C-R2N2 与 GNN 变体均呈现类似的 AIGC 可替换性，但 offline RL / 基于模型的 RL 或模仿学习可能呈现不同响应。

\subsection{部署实践指导}
\label{sec:guidance}

本文结果提供四点实践指导。其一，\textbf{容量规划比调度器选择更重要}：当 70B 请求溢出唯一云端服务器时（\S\ref{sec:workload}），所有调度器 \slo{} 都被压到 ${\approx}10\%$---\emph{先备算力，再优调度}。其二，\textbf{中等负载工况下简单负载均衡器足矣}：LeastLoaded 与 ShortestQueue 以零训练成本达到 RL \slo{} 的 $78\%$。其三，\textbf{当能效、尾延迟或分派时间预算重要时再上 RL}：RL 在全部 30 种配置下取得最低 \etok{}（负载下优势最高达 9\%），其尾 TTFT（P95 25.2~s vs.\ 28--41~s）适合交互式服务，且其亚毫秒级策略推理避免了 NSGA-II 类调度器在请求路径上每次分派 $\approx$70~ms 的搜索开销。其四，\textbf{对奖励工程的声明保持怀疑}：若无针对含预训练基线的联合消融，增益可能反映通用 RL 主干的提升，而非所声称的机制。
